% Necoyoad Architecture Blueprint v9 — The Banner Subsystem
\documentclass[11pt,a4paper]{article}
\usepackage[T1]{fontenc}
\usepackage[utf8]{inputenc}
\usepackage{lmodern}
\usepackage{microtype}
\usepackage{graphicx}
\usepackage{xcolor}
\usepackage{geometry}
\usepackage{amsmath}
\usepackage{amssymb}
\usepackage{tcolorbox}
\tcbuselibrary{skins,breakable,listings}
\usepackage{colortbl}
\usepackage{booktabs}
\usepackage{tabularx}
\usepackage{longtable}
\usepackage{array}
\usepackage{enumitem}
\usepackage{adjustbox}
\usepackage{caption}
\usepackage{fancyhdr}
\usepackage{titlesec}
\usepackage{pdflscape}
\usepackage{listings}
\usepackage{tikz}
\usetikzlibrary{arrows.meta,positioning,calc,shapes.geometric,fit,backgrounds,chains,decorations.pathreplacing,shapes.multipart}
\usepackage[colorlinks=true,linkcolor=blue!50!black,citecolor=blue!40!black,urlcolor=blue!60!black,bookmarks=true,bookmarksnumbered=true,unicode=true,pdftitle={Necoyoad - The Banner Subsystem (v9)},pdfauthor={Reverse-Engineered Architecture Report},pdfsubject={Software Architecture, Banners, Sliders, jQuery Plugins, Widget Composition},pdfkeywords={Necoyoad, OpenCart, PHP 8.1, Banner, Banner Item, Slider, jQuery Plugin, NivoSlider, Slick, ntPlugins, Widget Composition}]{hyperref}
\geometry{a4paper, top=2.4cm, bottom=2.4cm, left=2.4cm, right=2.4cm}
\setlength{\headheight}{14pt}
\pagestyle{fancy}
\fancyhf{}
\fancyhead[L]{\small\itshape Necoyoad v9 --- The Banner Subsystem}
\fancyfoot[C]{\small\thepage}
\renewcommand{\headrulewidth}{0.3pt}
\renewcommand{\footrulewidth}{0pt}
\titleformat{\section}{\Large\bfseries\color{black!85}}{\thesection}{0.6em}{}
\titleformat{\subsection}{\large\bfseries\color{black!80}}{\thesubsection}{0.6em}{}
\titleformat{\subsubsection}{\normalsize\bfseries\color{black!75}}{\thesubsubsection}{0.5em}{}
\widowpenalty=10000
\clubpenalty=10000
\emergencystretch=4em
\hbadness=10000
\hfuzz=2pt
\vfuzz=2pt
\sloppy
\definecolor{necoyoad-blue}{HTML}{1F4E79}
\definecolor{necoyoad-orange}{HTML}{B85C00}
\definecolor{necoyoad-green}{HTML}{2E6B33}
\definecolor{nodefill}{HTML}{F4F7FB}
\definecolor{nodefillalt}{HTML}{FBF4EE}
\definecolor{nodefillgreen}{HTML}{F2F8F2}
\definecolor{nodefillgray}{HTML}{F0F0F2}
\definecolor{phpline}{HTML}{F8F8F0}
\definecolor{phpcomment}{HTML}{8C8C8C}
\definecolor{phpkeyword}{HTML}{1F4E79}
\definecolor{phpstring}{HTML}{B85C00}
\lstdefinelanguage{necophp}{
  morekeywords={abstract,and,array,as,break,case,catch,class,clone,const,continue,declare,default,do,echo,else,elseif,empty,enddeclare,endfor,endforeach,endif,endswitch,endwhile,extends,final,finally,for,foreach,function,global,goto,if,implements,include,include_once,instanceof,insteadof,interface,namespace,new,or,print,private,protected,public,require,require_once,return,static,switch,throw,trait,try,use,var,while,xor,yield,true,false,null},
  sensitive=true,
  morecomment=[l]{//},
  morecomment=[s]{/*}{*/},
  morecomment=[l]{\#},
  morestring=[b]",
  morestring=[b]',
}
\lstset{
  language=necophp,
  basicstyle=\ttfamily\footnotesize,
  keywordstyle=\color{phpkeyword}\bfseries,
  commentstyle=\color{phpcomment}\itshape,
  stringstyle=\color{phpstring},
  backgroundcolor=\color{phpline},
  numbers=left,
  numberstyle=\tiny\color{phpcomment},
  numbersep=8pt,
  frame=single,
  rulecolor=\color{black!25},
  framerule=0.3pt,
  breaklines=true,
  breakatwhitespace=true,
  showstringspaces=false,
  captionpos=b,
  xleftmargin=14pt,
  xrightmargin=4pt,
  aboveskip=8pt,
  belowskip=8pt,
}
\newtcolorbox{findingbox}[1][]{
  enhanced, breakable,
  colback=nodefillalt, colframe=necoyoad-orange!70,
  boxrule=0.4pt, arc=2pt, left=8pt, right=8pt, top=6pt, bottom=6pt,
  fonttitle=\bfseries\small, coltitle=white,
  title=#1
}
\captionsetup{justification=raggedright, singlelinecheck=false, font=small, labelfont=bf}

\begin{document}
\thispagestyle{empty}
\begin{center}
{\Large\bfseries Necoyoad --- The Banner Subsystem (v9)}\\[0.4em]
{\small\itshape Banners, banner items, jQuery slider plugins,\\
and the 33-template rendering system}\\[1.2em]
\end{center}
\begin{abstract}
\noindent
This is the ninth volume of the Necoyoad architecture blueprint. v9 covers
the banner subsystem: how the admin creates banners with a
\texttt{jquery\_plugin} discriminator, how banner items (slides) carry
images, links, and localised descriptions, how the storefront banner
widget module loads the banner and dynamically selects the matching
slider template, how the \texttt{ntPlugins} JavaScript registry
initialises the jQuery slider, and how per-item widgets can be attached
to individual slides.

The banner subsystem is the most visually diverse part of Necoyoad: 33
slider/gallery templates in the \texttt{choroni/banner/} folder, each
backed by a jQuery plugin (NivoSlider, Slick, Camera, EisSlideshow,
Fancybox, etc.) with its own JS and CSS assets. The
\texttt{jquery\_plugin} column on the \texttt{banner} table is the
discriminator: the storefront controller reads it, loads the matching
\texttt{.tpl} file, enqueues the matching \texttt{slider.js} and
\texttt{slider.css}, and the template pushes a configuration object
into the \texttt{window.ntPlugins} registry, which
\texttt{theme.js::loadNTPlugins()} initialises on DOM ready.

v9 documents all of this with 8 PHP code listings, 3 TikZ diagrams,
the full 33-template inventory, and a banner rendering lifecycle trace.
The stance is unchanged: \textbf{descriptive, not prescriptive}.
\end{abstract}
\vspace{0.6em}
\noindent\textbf{Keywords:} Necoyoad, OpenCart, PHP 8.1, banner,
banner item, slider, jQuery plugin, NivoSlider, Slick, Camera,
\texttt{ntPlugins}, \texttt{loadNTPlugins}, widget composition,
polymorphic description, EAV property.
\newpage
\tableofcontents
\newpage

% ============================================================
\section{Introduction}\label{sec:intro}
% ============================================================

\subsection{Nine Volumes}

v1--v8 covered the whole platform, source verification, the rendering
pipeline, the marketing subsystem, multi-store/multi-language, the
admin back-office, the menu links, and the CMS posts/pages. v9 covers
the banner subsystem.

\subsection{Scope}

v9 covers:
\begin{itemize}[leftmargin=1.6em,itemsep=2pt]
  \item The \texttt{banner} and \texttt{banner\_item} tables.
  \item The admin banner controller (\texttt{ControllerContentBanner})
        and model (\texttt{ModelContentBanner}).
  \item The storefront banner widget module
        (\texttt{ControllerModuleBanner}).
  \item The storefront banner model (with store-scoped query and
        publish-window check).
  \item The \texttt{jquery\_plugin} discriminator and dynamic template
        selection.
  \item The 33 slider/gallery templates in
        \texttt{choroni/banner/}.
  \item The \texttt{ntPlugins} JavaScript registry and
        \texttt{loadNTPlugins()} initialiser.
  \item The per-item widget composition (widgets attached to
        individual banner items).
  \item The \texttt{carousel()} AJAX endpoint.
\end{itemize}

% ============================================================
\section{The Banner Schema}\label{sec:schema}
% ============================================================

\subsection{The \texttt{banner} Table}

\begin{lstlisting}[caption={The \texttt{banner} table schema.},label={lst:banner-schema}]
CREATE TABLE `banner` (
  `banner_id`          int(11) NOT NULL,
  `name`               varchar(250) NOT NULL,
  `jquery_plugin`      varchar(150) NOT NULL,  -- e.g. 'nivo-slider', 'slick', 'camera-v1.3.4'
  `params`             text NOT NULL,           -- serialised PHP slider config
  `publish_date_start` date NOT NULL,
  `publish_date_end`   date NOT NULL,
  `status`             int(1) NOT NULL,
  `date_added`         datetime NOT NULL,
  `date_modified`      datetime NOT NULL
);
\end{lstlisting}

The \texttt{jquery\_plugin} column is the key discriminator: it names
the jQuery slider plugin to use for rendering. The storefront controller
reads this column and uses it to select the matching template file
(\texttt{choroni/banner/<jquery\_plugin>.tpl}), the matching JS file
(\texttt{web/assets/js/sliders/<jquery\_plugin>/slider.js}), and the
matching CSS file
(\texttt{web/assets/css/sliders/<jquery\_plugin>/slider.css}).

The \texttt{params} column stores a serialised PHP array of slider
configuration (effect, speed, transition type, etc.). However, the
storefront templates hardcode the configuration in the
\texttt{ntPlugins.push()} call rather than reading from
\texttt{params} --- so \texttt{params} is effectively unused by the
storefront (it may be used by the admin form for persistence).

\subsection{The \texttt{banner\_item} Table}

\begin{lstlisting}[caption={The \texttt{banner\_item} table schema.},label={lst:banner-item-schema}]
CREATE TABLE `banner_item` (
  `banner_item_id` int(11) NOT NULL,
  `banner_id`      int(11) NOT NULL,
  `image`          varchar(250) NOT NULL,  -- filesystem path to the slide image
  `link`           varchar(250) NOT NULL,  -- URL the slide links to
  `sort_order`     int(11) NOT NULL,
  `status`         int(1) NOT NULL DEFAULT '1'
);
\end{lstlisting}

Each banner item is a single slide: an image, a link, a sort order, and
a status. The slide's localised title and description are stored in the
polymorphic \texttt{description} table with
\texttt{object\_type = 'banner\_item'} and
\texttt{object\_id = <banner\_item\_id>}. Per-item properties (like
offsetX, offsetY for positioning overlaid widgets) are stored in the
EAV \texttt{property} table with the same
\texttt{(object\_id, object\_type)} pair.

% ============================================================
\section{The Admin Banner Controller and Model}\label{sec:admin}
% ============================================================

\subsection{The Admin Controller}

\texttt{app/admin/controller/content/banner.php} (283 lines) extends
\texttt{ControllerAdmin} (v6) with standard declarative CRUD. The
\texttt{\$form\_vars} array declares the expected POST fields:

\begin{lstlisting}[caption={\texttt{ControllerContentBanner} --- declarative form vars (abridged).},label={lst:banner-form-vars}]
protected array $form_vars = [
    'banner_id'           => ['name' => 'banner_id', 'type' => 'number'],
    'name'                => ['name' => 'name', 'type' => 'string'],
    'jquery_plugin'       => ['name' => 'jquery_plugin', 'type' => 'string'],
    'publish_date_start'  => ['name' => 'publish_date_start', 'type' => 'date'],
    'publish_date_end'    => ['name' => 'publish_date_end', 'type' => 'date'],
    'banner_stores'       => ['name' => 'banner_stores', 'type' => 'array'],
    'banner_items'        => ['name' => 'banner_items', 'type' => 'array'],
    'banner_properties'   => ['name' => 'banner_properties', 'type' => 'array'],
    'stores'              => ['name' => 'stores', 'type' => 'array'],
];
\end{lstlisting}

The \texttt{banner\_items} array carries the slide data (image, link,
sort\_order, status, descriptions, properties) from the admin form.
The \texttt{banner\_stores} array carries the store assignments (for
\texttt{object\_to\_store}). Because the controller extends
\texttt{AdminController}, all CRUD is inherited.

\subsection{The Admin Model}

\texttt{app/admin/model/content/banner.php} (324 lines) declares
\texttt{\$table = "banner"}, \texttt{\$object\_type = "banner"}, and
\texttt{\$description\_object\_type = "banner\_item"}. The
\texttt{on('save')} hook manages items:

\begin{lstlisting}[caption={\texttt{ModelContentBanner} --- the on('save') hook for items.},label={lst:banner-model-save}]
public function init() {
    $this->on("save", function ($args) {
        $data = $args[0];
        if ($data['items']) {
            foreach ($data['items'] as $key => $item) {
                $item['banner_id'] = $data['id'];
                $item['sort_order'] = $key;
                $this->setItem($item);
            }
        }
    });

    $this->on("delete", function ($args) {
        $data = $args[0];
        $this->deleteItems($data['id']);
    });
}
\end{lstlisting}

The \texttt{setItem()} method handles both insert and update of a
banner item: it checks if \texttt{banner\_item\_id} exists, updates or
inserts accordingly, then saves the descriptions (localised title and
HTML description) and properties (EAV) for the item.

The \texttt{deleteItems()} method cascades: deletes
\texttt{description} rows, \texttt{property} rows, and
\texttt{banner\_item} rows for the given \texttt{banner\_id}.

The \texttt{getById()} method loads a banner with its items, stores,
and properties:

\begin{lstlisting}[caption={\texttt{ModelContentBanner::getById()} --- loads banner with items.},label={lst:banner-getbyid}]
public function getById($id) {
    $results = parent::getAll(array('banner_id'=>$id));
    $results[0]['banner_items'] = $this->getItems(array('banner_id'=>$id));
    $results[0]['banner_stores'] = $this->getStores($id);
    $results[0]['banner_properties'] = $this->getAllProperties($id);
    return $results[0];
}
\end{lstlisting}

% ============================================================
\section{The Storefront Banner Widget Module}\label{sec:storefront}
% ============================================================

\texttt{app/shop/controller/module/banner.php} (113 lines) is the
storefront widget module. It extends
\texttt{ControllerModuleModuleController} (v3) and registers a
\texttt{module:settings} filter in \texttt{init()}.

\subsection{The \texttt{module:settings} Filter}

The filter is the heart of the banner widget. It:

\begin{enumerate}[leftmargin=1.6em,itemsep=2pt]
  \item Loads the banner by \texttt{settings['banner\_id']} via
        \texttt{modelBanner->getById()}.
  \item Loads the banner items via
        \texttt{modelBanner->getItems()}.
  \item For each item, loads its localised descriptions and its
        per-item widgets (via a recursive
        \texttt{NecoWidget::getWidgets()} call with
        \texttt{object\_type = 'banner\_item'} and
        \texttt{object\_id = <banner\_item\_id>}).
  \item Enqueues the slider's JS file
        (\texttt{web/assets/js/sliders/<plugin>/slider.js}).
  \item Enqueues the slider's CSS file
        (\texttt{web/assets/css/sliders/<plugin>/slider.css}).
  \item Selects the template:
        \texttt{choroni/banner/<jquery\_plugin>.tpl}
        (or the active theme's version, or falls back to
        \texttt{nivo-slider.tpl}).
\end{enumerate}

\begin{lstlisting}[caption={\texttt{ControllerModuleBanner::init()} --- the module:settings filter (abridged).},label={lst:banner-widget}]
public function init() {
    $this->addFilter("module:settings", function ($data) {
        $settings = $data['settings'];
        if (isset($settings['banner_id'])) {
            $this->load->model('content/banner');
            $this->data['banner'] = $this->modelBanner->getById((int)$settings['banner_id']);

            if (isset($this->data['banner']['banner_id'])) {
                $items = $this->modelBanner->getItems($this->data['banner']['banner_id']);
                foreach ($items as $k => $item) {
                    // Load localised descriptions
                    $items[$k]['descriptions'] = $this->modelBanner->getDescriptions($item['banner_item_id']);

                    // Load per-item widgets (widgets attached to this slide)
                    $w = new NecoWidget($this->registry, $this->Route);
                    $params = array(
                        'store_id' => STORE_ID,
                        'landing_page' => 'all',
                        'object_type' => 'banner_item',
                        'object_id' => $item['banner_item_id'],
                        'position' => 'main'
                    );
                    $widgets = $w->getWidgets($params, !$this->user->getId());
                    foreach ($widgets as $v) {
                        $s = (array)unserialize($v['settings']);
                        $this->children[$v['name']] = $s['route'];
                        $this->widget[$v['name']] = $v;
                        $items[$k]['widgets'][$v['name']] = array(
                            'name' => $v['name'],
                            'offsetY' => $s['offsetY'],
                            'offsetX' => $s['offsetX']
                        );
                    }
                }
                $this->data['banner']['items'] = $items;
            }

            // Select template based on jquery_plugin
            if (!empty($this->data['banner']['jquery_plugin'])) {
                $plugin = $this->data['banner']['jquery_plugin'];
                // Enqueue JS
                if (file_exists(DIR_JS . 'sliders/' . $plugin . '/slider.js')) {
                    $this->javascripts[$plugin] = HTTP_JS . 'sliders/' . $plugin . '/slider.js';
                }
                // Enqueue CSS
                if (file_exists(DIR_CSS . 'sliders/' . $plugin . '/slider.css')) {
                    $this->styles[] = ['href' => HTTP_CSS . 'sliders/' . $plugin . '/slider.css', 'media' => 'all'];
                }
                // Select template
                if (file_exists(DIR_TEMPLATE . $this->config->get('config_template') . '/banner/' . $plugin . '.tpl')) {
                    $this->template = $this->config->get('config_template') . '/banner/' . $plugin . '.tpl';
                } elseif (file_exists(DIR_TEMPLATE . 'choroni/banner/' . $plugin . '.tpl')) {
                    $this->template = 'choroni/banner/' . $plugin . '.tpl';
                } else {
                    $this->template = 'choroni/banner/nivo-slider.tpl';
                }
            }
        }
        return $data;
    });
}
\end{lstlisting}

\subsection{The Per-Item Widget Composition}

\begin{findingbox}[Banners support per-item widget composition]
Each banner item (slide) can have its own set of widgets, loaded via a
recursive \texttt{NecoWidget::getWidgets()} call with
\texttt{object\_type = 'banner\_item'} and
\texttt{object\_id = <banner\_item\_id>}. These widgets are registered
as children of the banner widget controller (merged into
\texttt{\$this->children}), and their \texttt{offsetX}/\texttt{offsetY}
settings are stored in the \texttt{\$items[\$k]['widgets']} array for
the template to position them as overlays on the slide.

This means a single banner slide can have overlaid widgets (e.g.\ a
``product overview'' widget positioned at \texttt{offsetX=100,
offsetY=50} over the slide image) --- the same widget composition
mechanism that drives page-level composition (v3, v8) is applied at
the slide level.
\end{findingbox}

\subsection{The \texttt{carousel()} AJAX Endpoint}

The banner module also provides a \texttt{carousel()} method that
returns banner items as JSON (with thumbnail resizing via
\texttt{NTImage::resizeAndSave()}). This is used by JavaScript carousels
that load their slides via AJAX rather than server-side rendering:

\begin{lstlisting}[caption={\texttt{ControllerModuleBanner::carousel()} --- the AJAX endpoint.},label={lst:banner-carousel}]
public function carousel() {
    $this->load->model('content/banner');
    $this->load->auto('image');
    $this->load->auto('json');

    $banner = $this->modelBanner->getById($this->request->getQuery('banner_id'));
    $json['results'] = $banner['items'];

    $width = isset($_GET['width']) ? $_GET['width'] : 80;
    $height = isset($_GET['height']) ? $_GET['height'] : 80;
    foreach ($json['results'] as $k => $v) {
        $json['results'][$k]['thumb'] = NTImage::resizeAndSave($v['image'], $width, $height);
        $json['results'][$k]['title'] = $v['title'];
        $json['results'][$k]['link'] = $v['link'];
    }

    $this->response->setOutput(Json::encode($json), $this->config->get('config_compression'));
}
\end{lstlisting}

% ============================================================
\section{The Storefront Banner Model}\label{sec:storefront-model}
% ============================================================

\texttt{app/shop/model/content/banner.php} (63 lines) provides
storefront read access to banners with store-scoping and
publish-window checks:

\begin{lstlisting}[caption={Storefront \texttt{ModelContentBanner::getById()} --- store-scoped with publish window.},label={lst:storefront-banner-model}]
public function getById($banner_id) {
    $query = $this->db->query("SELECT DISTINCT * FROM banner b
        LEFT JOIN object_to_store b2s ON
            (b.banner_id = b2s.object_id AND b2s.object_type = 'banner')
        WHERE b.banner_id = '" . (int)$banner_id . "'
        AND b.publish_date_start <= NOW()
        AND (b.publish_date_end >= NOW() OR b.publish_date_end = '0000-00-00')
        AND b.status = '1'
        AND b2s.store_id = '" . (int)STORE_ID . "'");

    $return = [];
    if ($query->num_rows) {
        $return = array_merge($query->row, array('items' => $this->getItems($banner_id)));
    }
    return $return;
}

public function getItems($banner_id) {
    $sql = "SELECT * FROM banner_item bi
        WHERE bi.banner_id = '" . (int)$banner_id . "'
        AND bi.status = '1'";
    return $this->db->query($sql)->rows;
}
\end{lstlisting}

The query checks:
\begin{itemize}[leftmargin=1.6em,itemsep=2pt]
  \item \texttt{publish\_date\_start <= NOW()} --- the banner's
        publish window has started.
  \item \texttt{publish\_date\_end >= NOW() OR publish\_date\_end =
        '0000-00-00'} --- the banner's publish window has not ended
        (or has no end date).
  \item \texttt{status = '1'} --- the banner is enabled.
  \item \texttt{b2s.store\_id = STORE\_ID} --- the banner is assigned
        to the current store (via the \texttt{object\_to\_store}
        junction, v5).
\end{itemize}

The items query only returns items with \texttt{status = '1'}
(enabled slides).

% ============================================================
\section{The 33 Slider/Gallery Templates}\label{sec:templates}
% ============================================================

The \texttt{choroni/banner/} folder contains 33 templates, each backed
by a jQuery plugin:

\begin{longtable}{>{\raggedright\arraybackslash}p{5.0cm}>{\raggedright\arraybackslash}p{8.0cm}}
\caption{The 33 banner templates.}\label{tab:banner-templates}\\
\toprule
\textbf{Template} & \textbf{Slider/gallery type} \\
\midrule
\endfirsthead
\toprule
\textbf{Template} & \textbf{Slider/gallery type} \\
\midrule
\endhead
\bottomrule
\endfoot
\texttt{nivo-slider.tpl}         & NivoSlider (the default fallback). \\
\texttt{nivo-slider-v3.1.tpl}    & NivoSlider v3.1. \\
\texttt{slick.tpl}               & Slick carousel. \\
\texttt{camera-v1.3.4.tpl}       & Camera (jQuery slideshow) v1.3.4. \\
\texttt{eislideshow.tpl}         & EIS Slideshow. \\
\texttt{evolution-v1.1.5.tpl}    & Evolution slider v1.1.5. \\
\texttt{layer-slider-v0.0.1.tpl} & Layer slider v0.0.1. \\
\texttt{parallax-content-slider.tpl} & Parallax content slider. \\
\texttt{horizontal-parallax.tpl} & Horizontal parallax slider. \\
\texttt{slicebox-v1.1.0.tpl}     & Slicebox 3D slider v1.1.0. \\
\texttt{jssor-vertical-nav.tpl}  & Jssor slider with vertical navigation. \\
\texttt{necoslider.tpl}          & Necoyoad's own slider (custom). \\
\texttt{carousel.tpl}            & Simple carousel. \\
\texttt{horizontal.tpl}          & Horizontal slider. \\
\texttt{vertical.tpl}            & Vertical slider. \\
\texttt{fancybox-gallery.tpl}    & Fancybox lightbox gallery. \\
\texttt{fancybox-grid.tpl}       & Fancybox grid gallery. \\
\texttt{grid-gallery.tpl}        & CSS grid gallery. \\
\texttt{horizontal-hover\_effect-01} through \texttt{-15} & 15 horizontal hover-effect gallery variants. \\
\bottomrule
\end{longtable}

\subsection{Template Structure}

Each slider template follows the same pattern:

\begin{lstlisting}[caption={The universal slider template pattern (abridged from \texttt{nivo-slider.tpl}).},label={lst:slider-pattern}]
<?php $tpl = is_dir(DIR_TEMPLATE. $this->config->get('config_template') ."/shared")
    ? $this->config->get('config_template') : "choroni"; ?>
<li id="<?php echo $widgetName; ?>" class="banner nivo nt-editable"
    data-banner="nivoSlider">

<?php if (count($banner['items'])) { ?>
    <div class="content">
        <!-- The slider container -->
        <div id="slider" class="nivoSlider">
            <?php foreach ($banner['items'] as $item) { ?>
                <img src="<?php echo HTTP_IMAGE . $item['image']; ?>"
                     alt="<?php echo $item['descriptions'][$Config->get('config_language_id')]['title']; ?>"
                     title="<?php echo $item['descriptions'][$Config->get('config_language_id')]['title']; ?>" />
            <?php } ?>
        </div>
    </div>

    <!-- Push slider config to the ntPlugins registry -->
    <script>
    $(function(){
        ntPlugins = window.ntPlugins || [];
        ntPlugins.push({
            id: "<?php echo $widgetName; ?> .nivoSlider",
            config: {
                effect: 'random',
                slices: 12,
                animSpeed: 300,
                pauseTime: 6000,
                // ... more slider-specific options
            },
            plugin: 'nivoSlider'
        });
        window.ntPlugins = ntPlugins;

        if (typeof loadNTPlugins !== 'undefined' && typeof loadNTPlugins === 'function') {
            loadNTPlugins();
        }
    });
    </script>
<?php } ?>
</li>
\end{lstlisting}

The template:
\begin{enumerate}[leftmargin=1.6em,itemsep=2pt]
  \item Iterates \texttt{\$banner['items']} and emits the slider's
        HTML structure (e.g.\ \texttt{<img>} tags inside a
        \texttt{<div class="nivoSlider">}).
  \item Uses \texttt{\$Config->get('config\_language\_id')} to select
        the localised title/description for each slide (from the
        polymorphic \texttt{description} table).
  \item Uses \texttt{\$Image->resizeAndSave()} for thumbnail
        generation.
  \item Pushes a configuration object into
        \texttt{window.ntPlugins} with the slider's jQuery plugin name
        and options.
  \item Calls \texttt{loadNTPlugins()} to initialise the plugin.
\end{enumerate}

\subsection{The \texttt{ntPlugins} Registry and \texttt{loadNTPlugins()}}

The \texttt{ntPlugins} registry (v3, \S2.5) is a JavaScript array that
collects slider initialisation configurations. Each template pushes an
entry with:
\begin{itemize}[leftmargin=1.6em,itemsep=2pt]
  \item \texttt{id} --- the CSS selector for the slider container.
  \item \texttt{config} --- the slider-specific options (effect,
        speed, transition, etc.).
  \item \texttt{plugin} --- the jQuery plugin method name (e.g.\
        \texttt{'nivoSlider'}, \texttt{'slick'}).
\end{itemize}

\texttt{theme.js::loadNTPlugins()} iterates the array and, for each
entry, calls \texttt{\$(item.id)[item.plugin](item.config)} --- i.e.\
it applies the jQuery plugin to the slider container with the
configuration. After processing, each entry is set to \texttt{null}
to prevent re-initialisation.

This decoupled registry pattern allows multiple sliders on the same
page (each with its own plugin and config) without the templates
needing to know about each other.

% ============================================================
\section{The Banner Rendering Lifecycle}\label{sec:trace}
% ============================================================

\subsection{Admin Creates a Banner}

\begin{enumerate}[leftmargin=1.6em,itemsep=2pt]
  \item Admin navigates to \texttt{?r=content/banner/insert}.
  \item \texttt{ControllerContentBanner::insert()} (inherited from
        AdminController) calls \texttt{upsert()}.
  \item \texttt{upsert()} validates the form, processes POST, calls
        \texttt{modelBanner->add(\$data)}.
  \item \texttt{ModelContentBanner::add()} inserts the \texttt{banner}
        row (name, jquery\_plugin, params, publish dates, status).
  \item The \texttt{on('save')} hook fires: iterates
        \texttt{\$data['items']}, calls \texttt{setItem()} per slide,
        which inserts \texttt{banner\_item} rows + \texttt{description}
        rows (localised titles) + \texttt{property} rows (per-item
        settings).
  \item The \texttt{object\_to\_store} rows are written (via the
        \texttt{relations = ["stores"]} declaration) to scope the
        banner to stores.
  \item Admin activity logged.
\end{enumerate}

\subsection{Storefront Renders the Banner Widget}

\begin{enumerate}[leftmargin=1.6em,itemsep=2pt]
  \item A page loads with a ``banner'' widget configured with
        \texttt{banner\_id = 5}.
  \item \texttt{ControllerModuleBanner::init()} registers the
        \texttt{module:settings} filter.
  \item \texttt{ControllerModuleModuleController::index()} calls the
        filter.
  \item The filter:
        \begin{enumerate}[itemsep=2pt]
          \item Calls \texttt{modelBanner->getById(5)} ---
                \texttt{SELECT * FROM banner b LEFT JOIN
                object\_to\_store b2s ON (...) WHERE b.banner\_id = 5
                AND b.publish\_date\_start <= NOW() AND ... AND
                b2s.store\_id = STORE\_ID}.
          \item Calls \texttt{modelBanner->getItems(5)} ---
                \texttt{SELECT * FROM banner\_item WHERE banner\_id =
                5 AND status = '1'}. Returns the slides.
          \item For each slide, loads descriptions
                (\texttt{getDescriptions(banner\_item\_id)}) and
                per-item widgets
                (\texttt{NecoWidget::getWidgets(...)} with
                \texttt{object\_type = 'banner\_item'}).
          \item Reads \texttt{banner['jquery\_plugin']} (e.g.\
                \texttt{'nivo-slider'}).
          \item Enqueues \texttt{sliders/nivo-slider/slider.js} and
                \texttt{sliders/nivo-slider/slider.css}.
          \item Sets template to
                \texttt{choroni/banner/nivo-slider.tpl}.
        \end{enumerate}
  \item \texttt{ControllerModuleModuleController::index()} calls
        \texttt{render()}.
  \item \texttt{fetch('choroni/banner/nivo-slider.tpl')}:
        \begin{enumerate}[itemsep=2pt]
          \item Injects magic services (\$Config, \$Language, \$l,
                \$Request, \$Url, \$Image).
          \item Extracts \texttt{\$this->data} into local scope
                (\$banner, \$widgetName, \$settings, etc.).
          \item Requires the template. The template:
                \begin{itemize}[itemsep=2pt]
                  \item Iterates \texttt{\$banner['items']}, emits
                        \texttt{<img>} tags with localised titles.
                  \item Pushes a config object into
                        \texttt{window.ntPlugins}.
                  \item Calls \texttt{loadNTPlugins()}.
                \end{itemize}
          \item Substitutes \texttt{\{\%<per-item-widget-name>\%\}}
                placeholders with the per-item widget HTML.
          \item Returns the HTML.
        \end{enumerate}
  \item The HTML is returned via the \texttt{module:settings} filter
        and stored in \texttt{\$this->data['<widget\_name>\_code']}.
  \item The parent page's \texttt{fetch()} substitutes the
        \texttt{\{\%<widget\_name>\%\}} placeholder with the banner
        widget's HTML.
  \item On the client side, \texttt{theme.js} runs
        \texttt{loadNTPlugins()} on DOM ready, which applies the
        NivoSlider plugin to the slider container.
\end{enumerate}

% ============================================================
\section{Cross-Cutting Findings}\label{sec:findings}
% ============================================================

\subsection{The Banner Subsystem is a Plugin-Aware Widget}

The banner subsystem is the only widget module that dynamically
selects its template, JS, and CSS based on a database column
(\texttt{jquery\_plugin}). All other widget modules have a fixed
template (selected by the admin's \texttt{settings['view']} field).
The banner widget reads the \texttt{jquery\_plugin} column from the
\texttt{banner} table and uses it to look up
\texttt{choroni/banner/<plugin>.tpl},
\texttt{sliders/<plugin>/slider.js}, and
\texttt{sliders/<plugin>/slider.css}.

\subsection{The \texttt{ntPlugins} Registry is the JS Initialisation Seam}

The \texttt{ntPlugins} registry is the JavaScript initialisation seam
for all slider templates. Each template pushes a
\texttt{\{id, config, plugin\}} entry, and
\texttt{loadNTPlugins()} applies the jQuery plugin to the DOM element.
This is a decoupled, deferred-initialisation pattern: templates don't
need to know whether the jQuery plugin is loaded yet (the
\texttt{loadNTPlugins()} function checks
\texttt{\$[item.plugin] || \$.fn[item.plugin]} before applying).

\subsection{Per-Item Widget Composition is Unique to Banners}

No other widget module supports per-item widget composition. The
banner widget's \texttt{module:settings} filter runs a recursive
\texttt{NecoWidget::getWidgets()} call for each slide, loading
widgets scoped to \texttt{object\_type = 'banner\_item'} and
\texttt{object\_id = <banner\_item\_id>}. These widgets are registered
as children of the banner controller and their HTML is substituted
into the slide template. The \texttt{offsetX}/\texttt{offsetY}
settings allow the admin to position widgets as overlays on the
slide image.

\subsection{The \texttt{params} Column is Effectively Unused}

The \texttt{banner.params} column stores a serialised PHP array of
slider configuration, but the storefront templates hardcode the
configuration in the \texttt{ntPlugins.push()} call. The
\texttt{params} column may be used by the admin form for persistence
(saving the slider configuration), but it is not read by the
storefront. Each template has its own hardcoded config, which means
changing the slider speed or effect requires editing the template
file, not the admin UI.

\subsection{The \texttt{banner\_items} Form Array Carries Everything}

The admin form submits all slide data (images, links, descriptions,
properties) in a single \texttt{banner\_items} POST array. The
model's \texttt{on('save')} hook iterates this array and calls
\texttt{setItem()} per slide, which handles insert-or-update,
descriptions, and properties. This is a bulk-save pattern: all slides
are saved in one form submission, not via individual AJAX calls.

% ============================================================
\section{Closing Observations}\label{sec:observations}
% ============================================================

\subsection{The Banner Subsystem is the Most Visually Diverse Part}

With 33 templates backed by 14+ jQuery plugins, the banner subsystem
is the most visually diverse part of Necoyoad. The
\texttt{jquery\_plugin} discriminator allows the admin to choose a
different slider or gallery for each banner, and the storefront
dynamically loads the matching template, JS, and CSS. This is a
plugin-aware widget architecture that rivals dedicated slider
libraries.

\subsection{The Per-Item Widget Composition is a Unique Feature}

The ability to attach widgets to individual slides --- not just to
the banner as a whole --- is a unique feature. It allows the admin
to create rich, interactive slideshows where each slide has
overlaid content (product overviews, text blocks, contact forms)
positioned at specific coordinates. This is the same widget
composition mechanism (v3, v8) applied at the slide level, using
the same \texttt{NecoWidget::getWidgets()} query with
\texttt{object\_type = 'banner\_item'}.

\subsection{Final Note}

v9 confirms that the Necoyoad banner subsystem is a plugin-aware,
widget-composed slideshow system. The \texttt{jquery\_plugin}
discriminator, the 33-template library, the \texttt{ntPlugins}
registry, and the per-item widget composition together form a
flexible visual content system that can render everything from a
simple NivoSlider to a complex parallax gallery with overlaid
widgets.

\appendix
\section{Glossary}\label{app:glossary}
\begin{table}[htbp]
\centering
\caption{Glossary of banner subsystem terms.}\label{tab:glossary-v9}
\footnotesize\sloppy
\begin{tabularx}{\columnwidth}{>{\raggedright\arraybackslash}p{3.5cm}X}
\toprule
\textbf{Term} & \textbf{Meaning} \\
\midrule
Banner & A named collection of slides, with a \texttt{jquery\_plugin} discriminator. \\
Banner item & A single slide: image, link, sort\_order, status, localised descriptions, per-item properties. \\
\texttt{jquery\_plugin} & The discriminator column that names the jQuery slider plugin to use. Drives template, JS, and CSS selection. \\
\texttt{ntPlugins} & A JavaScript array registry. Each slider template pushes a \texttt{\{id, config, plugin\}} entry. \texttt{loadNTPlugins()} applies the plugin on DOM ready. \\
Per-item widget & A widget attached to an individual slide via \texttt{object\_type = 'banner\_item'}. Positioned with \texttt{offsetX}/\texttt{offsetY}. \\
\texttt{carousel()} & An AJAX endpoint that returns banner items as JSON with thumbnails. \\
\bottomrule
\end{tabularx}
\end{table}
\end{document}
